\documentclass[12pt, a4paper]{article}

\usepackage[utf8]{inputenc}
\usepackage[portuguese]{babel}
\usepackage[T1]{fontenc}
\usepackage{geometry}
\usepackage{graphicx}
\usepackage{listings}
\usepackage{xcolor}
\usepackage{booktabs}
\usepackage{float}
\usepackage{amsmath}
\usepackage{hyperref}
\usepackage{setspace}
\usepackage{array}
\usepackage{tabularx}

% Configuração das margens
\geometry{top=3cm, bottom=2cm, left=3cm, right=2cm}

% Configuração para códigos
\definecolor{codegreen}{rgb}{0,0.6,0}
\definecolor{codegray}{rgb}{0.5,0.5,0.5}
\definecolor{codepurple}{rgb}{0.58,0,0.82}
\definecolor{backcolour}{rgb}{0.95,0.95,0.92}

\lstdefinestyle{mystyle}{
    backgroundcolor=\color{backcolour},
    commentstyle=\color{codegreen},
    keywordstyle=\color{magenta},
    numberstyle=\tiny\color{codegray},
    stringstyle=\color{codepurple},
    basicstyle=\ttfamily\footnotesize,
    breakatwhitespace=false,
    breaklines=true,
    captionpos=b,
    keepspaces=true,
    numbers=left,
    numbersep=5pt,
    showspaces=false,
    showstringspaces=false,
    showtabs=false,
    tabsize=2
}

\lstset{style=mystyle}

\begin{document}

% Capa
\begin{titlepage}
  \centering
  {\LARGE \textbf{Universidade de São Paulo (USP)}}\\
  Escola de Artes, Ciências e Humanidades\\
  \vspace{2cm}
  {\Huge \textbf{EP - Sistema Distribuído PeerSpot}\\}
  \vspace{1.5cm}
  {\Large Disciplina: Desenvolvimento de Sistemas de Informação Distribuídos}\\
  Professor: Renan Cerqueira Afonso Alves\\
  \vspace{1.5cm}
  {\Large Alunos:}\\
  \begin{spacing}{1.2}
    Bruno Hideo Ioneda --- 15573619\\
    Guilherme Samuel Lemos Segura --- 15575611\\
    Higor Ranel Viani Lopes --- 15552946\\
    João de Melo Fantini --- 15462550\\
  \end{spacing}
  \vfill
  {\Large São Paulo, março de 2026}
\end{titlepage}

\newpage
\tableofcontents
\newpage

% ============================================================
% AULA 1 --- ARQUITETURA
% ============================================================

% ============================================================
\section{Dos Modelos de Arquitetura Estudados, algum se encaixa?}\label{sec:modelos}
% ============================================================

O modelo que melhor se encaixa para o projeto é o \textbf{P2P Não Estruturado Hierárquico com super peers}, descrito por Tanenbaum no Cap. 2.4.3. Nessa arquitetura, os nós da rede se dividem em dois tipos de \textbf{peers}:
\begin{itemize}
    \item \textbf{peers comuns (weak peers)}: ligados como clientes a um super peer
    \item \textbf{super peers}: mantêm um índice dos arquivos disponíveis e organizam o acesso.
\end{itemize}

Ao longo deste documento utilizamos o termo \textbf{tracker} no lugar de ``super peer'', em referência ao papel análogo desempenhado pelo tracker do BitTorrent (mencionado por Tanembaum no  Cap. 2.4.4), que pode ser entendido como um servidor que mantém uma conta atualizada de quais peers têm quais chunks de quais arquivos.

No PeerSpot, os trackers acumulam, portanto, duas funções: 
\begin{itemize}
    \item o papel de \textbf{super peer}
    \item o papel de \textbf{tracker}
\end{itemize}


\subsection{Por que não cliente-servidor puro?}

A arquitetura cliente-servidor ancora todo o sistema em um único componente central, que precisa ser robusto o suficiente para atender a todos os usuários simultaneamente. Em um serviço de compartilhamento de arquivos de áudio, o servidor central se tornaria rapidamente um gargalo de largura de banda e um ponto único de falha.

\subsection{Por que não P2P não estruturado puro?}

Uma alternativa considerada foi o P2P não estruturado puro, em que cada nó conhece apenas alguns vizinhos e as buscas são feitas por flooding ou random walks. O problema é que o flooding gera muitas mensagens: como cada nó repassa a busca para todos os vizinhos, o tráfego cresce rápido conforme a rede aumenta, o que não escala bem. Já os random walks diminuem esse tráfego, mas em compensação deixam a busca mais lenta.

Como nosso sistema é voltado para compartilhamento de música, em que o usuário espera encontrar o que procura rapidamente, nenhuma das duas opções nos pareceu adequada. Por isso optamos pela abordagem com super peers, que é justamente o que o livro recomenda (Cap. 2.4.3) quando trata da escalabilidade da busca em redes não estruturadas.

\subsection{Por que não P2P estruturado com DHT?}\label{sec:dht}
Outra opção avaliada foi o P2P estruturado com DHT, que oferece buscas rápidas sem componentes centrais. O gargalo da arquitetura seria o trabalho extra para manter as tabelas de roteamento e redistribuir as chaves sempre que um nó entra ou sai da rede. 

Tanembaum cita no Cap. 2.4.4 uma variante do BitTorrent que troca os trackers por uma DHT auxiliar, mas descartamos essa ideia pelo mesmo motivo: a complexidade adicional não compensa para o escopo do trabalho.

\subsection{Por que P2P com super peers (trackers)?}

A arquitetura com super peers traz os seguintes benefícios:

Os trackers mantêm um \textbf{índice replicado} entre vários nós (sem ponto único de falha), enquanto a transferência dos arquivos ocorre diretamente entre os peers, sem passar por nenhum intermediário. Isso traz:

\begin{itemize}
    \item \textbf{Escalabilidade:} a carga de armazenamento e transferência é distribuída entre os peers. Adicionar usuários aumenta a capacidade da rede.
    \item \textbf{Tolerância a falhas:} como os trackers sincronizam seus índices entre si por \textit{flooding} sobre TCP, a queda de um tracker não torna os arquivos inacessíveis.
    \item \textbf{Eficiência de busca:} weak peers não precisam manter tabelas de roteamento complexas, delegam a busca ao tracker ao qual estão conectados.
    \item \textbf{Transparência de localização}: peers acessam arquivos por nome legível ou hash, sem saber qual peer físico vai servir a transferência.
\end{itemize}

\subsection{Componentes e Responsabilidades}\label{sec:componentes}

Definida a arquitetura, a seguir descrevemos os dois tipos de componente e como o conjunto de trackers se organiza.

\paragraph{Peer}
\begin{itemize}
    \item Nó básico da rede.
    \item Armazena chunks de arquivos de áudio localmente.
    \item Mantém uma conexão duradoura com seu tracker.
    \item Ao fazer upload de uma música, notifica o tracker com \texttt{nome + hash}.
    \item Para baixar, consulta o tracker, recebe a lista de peers que possuem os chunks e realiza a transferência diretamente via TCP.
\end{itemize}

\paragraph{Tracker}
\begin{itemize}
    \item Agrega um cluster de peers comuns. Mantém cinco estruturas de índice em memória:
    \begin{itemize}
        \item \texttt{nome\_da\_musica} $\rightarrow$ \texttt{conjunto de hashes} (homônimos são suportados)
        \item \texttt{hash} $\rightarrow$ \texttt{metadados} (\texttt{nome}, \texttt{tamanho}, \texttt{n\_chunks})
        \item \texttt{hash} $\rightarrow$ \texttt{lista de peers que possuem o arquivo}
        \item \texttt{nome\_peer} $\rightarrow$ \texttt{(ip, porta, last\_seed\_ts)}
        \item \texttt{tombstones} (marcações de remoção com \textit{timestamp}, retidas por 10 minutos)
    \end{itemize}
    \item Responde a buscas dos peers do seu cluster e roteia para outros trackers quando necessário.
    \item Sincroniza suas tabelas com os demais trackers.
\end{itemize}

\paragraph{Organização dos trackers}
\begin{itemize}
    \item Os trackers formam uma rede totalmente conectada: cada tracker mantém uma lista de todos os outros trackers ativos.
    \item  Não há hierarquia entre eles, todos são equivalentes e mantêm cópias sincronizadas do índice global.
\end{itemize}

% ============================================================
\section{Arquitetura de Software Interna}\label{sec:arq-interna}
% ============================================================

Cada componente possui uma organização interna distinta que serão detalhadas nas seções seguintes.

\subsection{Servidor de Metadados / tracker}

Adota \textbf{arquitetura em camadas}:

\begin{itemize}
    \item \textbf{Camada de API REST:} recebe requisições dos peers (\texttt{REGISTER\_FILE}, \texttt{SEARCH\_FILE}) e dos demais trackers em caso de roteamento.
    \item \textbf{Camada de lógica de negócio:} processa autenticação, gerenciamento de playlists, lookup nas tabelas de índice e roteamento de buscas.
    \item \textbf{Camada de persistência:} banco de dados relacional para dados duráveis (usuários, playlists). As tabelas de índice (como \texttt{nome→hash} e \texttt{hash→peers}) são mantidas em memória, pois mudam com frequência e devem ser reconstruídas a partir dos peers ativos após uma reinicialização.
    \item \textbf{Módulo de sincronização:} responsável por enviar e receber mensagens \texttt{SYNC\_TABLE} via TCP unicast (uma conexão paralela para cada tracker conhecido) para manter os índices sincronizados entre todos os trackers, e por gerenciar a entrada/saída de trackers através de \texttt{TRACKER\_REJOIN}, \texttt{FULL\_SYNC} e \texttt{TRACKER\_ANNOUNCE}.
\end{itemize}

\subsection{Peer / Cliente}

Adota \textbf{arquitetura P2P}: não há distinção de papel fixo, o mesmo nó baixa e serve chunks simultaneamente.

\begin{itemize}
    \item \textbf{Interface:} Mantém a CLI responsiva as ações do usuário
    \item \textbf{Módulo P2P:} envia e recebe chunks de arquivos diretamente com outros peers via TCP.
    \item \textbf{Gerenciador de chunks:} controla quais partes de cada arquivo foram obtidas e quais ainda faltam. Antes de iniciar um download, consulta cada peer fonte com \texttt{CHUNK\_LIST\_REQUEST} para descobrir os chunks disponíveis em cada uma, e então coordena downloads paralelos distribuindo as requisições entre as múltiplas fontes. E se a conexão for interrompida retoma o download considerando os chunks que já foram baixados.
    \item \textbf{Armazenamento local:} Arquivos de áudio, e quando solicitado segmenta o arquivo em chunks para propagar.
\end{itemize}

% ============================================================
\section{Como o Sistema Será Testado}\label{sec:testes}
% ============================================================

A seguir, definimos como o sistema será testado, considerando o escopo do projeto.

\subsection{Testes Funcionais Locais}
Simular múltiplos peers na mesma máquina utilizando processos / terminais diferentes. O objetivo é validar as funcionalidades básicas, upload, registro no tracker, busca, download e reprodução, em um ambiente controlado, antes de introduzir variáveis de rede real.


\subsection{Testes de Tolerância a Falhas}
Derrubar tracker e peers deliberadamente para verificar:
\begin{itemize}
    \item Se os arquivos continuam acessíveis por meio de outros peers que possuem os chunks.
    \item Se a sincronização entre tracker se mantém após a saída de um nó da rede de trackers.
    \item Se peers órfãos reconectam corretamente ao tracker de backup após falha do tracker principal.
\end{itemize}

% ============================================================
\section{Middleware}\label{sec:middleware}
% ============================================================

O Tanenbaum (Cap.~2.2) define \textit{middleware} como a camada que abstrai a heterogeneidade da rede, permitindo que componentes distribuídos se comuniquem sem conhecer os detalhes uns dos outros. No PeerSpot, o middleware se manifesta entre os peers e entre um peer e seu tracker correspondente.

\subsection{Entre peer e tracker}
Um framework de comunicação padronizado, REST sobre HTTP, abstrai os detalhes de rede e permite que peers em plataformas diferentes (Windows, Linux, macOS) se comuniquem com o tracker de forma uniforme. Bibliotecas como \texttt{Flask} ou \texttt{FastAPI} (Python) atuam como middleware nessa camada.

\subsection{Entre peers}
O protocolo de troca de chunks sobre TCP padroniza como dois peers negociam quais partes de um arquivo um tem e o outro precisa. A biblioteca padrão de \textit{sockets} POSIX (módulo \texttt{socket} do Python) combinada com \texttt{threading} atua como middleware aqui, isolando o módulo P2P dos detalhes de baixo nível da comunicação.

O uso de middleware se justifica principalmente pela \textbf{heterogeneidade} do sistema: como os peers podem rodar em sistemas operacionais diferentes, o middleware é o que resolve o problema de fazer todos conversarem entre si. Além disso, ao adotar REST/JSON e sockets POSIX padrão, usamos formatos e protocolos abertos e padronizados, o que permite que implementações diferentes do peer ou do tracker consigam funcionar juntas sem problema.

% ============================================================
\section{Diagrama de Arquitetura}\label{sec:diag-arq}
% ============================================================
\begin{figure}[H]
    \centering
    \includegraphics[width=1\linewidth]{desenho_arquitetura.jpeg}
    \caption{Diagrama de Arquitetura}
    \label{fig:arquitetura}
\end{figure}

% ============================================================
% AULA 2 --- COMUNICAÇÃO
% ============================================================

% ============================================================
\section{Comunicação}\label{sec:comunicacao}
% ============================================================

\subsection{Síncrona ou Assíncrona?}

A escolha entre comunicação síncrona e assíncrona vai depender da operação:

\begin{itemize}
    \item \textbf{Síncrona}, operações em que o peer precisa da resposta para continuar: busca de música (\texttt{SEARCH\_FILE}) e registro de arquivo (\texttt{REGISTER\_FILE}). O peer aguarda a resposta do tracker antes de prosseguir.
    
    \item \textbf{Assíncrona}, sincronização entre trackers via flooding sobre TCP. O tracker que atualizou seu índice não aguarda confirmação dos demais: dispara o \texttt{SYNC\_TABLE} para cada tracker conhecido em threads paralelas e continua operando. Isso evita que a indisponibilidade temporária de um tracker bloqueie os demais. A convergência é garantida pela reconciliação \textit{anti-entropy} periódica entre trackers (Seção~\ref{sec:rep-protocolo}).
\end{itemize}

\subsection{Protocolos de Transporte}

\begin{center}
\begin{tabularx}{\textwidth}{lXX}
\toprule
\textbf{Comunicação} & \textbf{Protocolo} & \textbf{Justificativa} \\
\midrule
Peer $\leftrightarrow$ tracker & REST sobre TCP (HTTP) & Simplicidade, padronização, suporte em qualquer linguagem \\
tracker $\leftrightarrow$ tracker (sincronização) & TCP (flooding) & Entrega ordenada e confiável, detecção natural de tracker falho via timeout/erro de conexão \\
Peer $\leftrightarrow$ Peer (transferência de chunks) & TCP direto & Garante entrega ordenada e sem perdas para dados binários \\
\bottomrule
\end{tabularx}
\end{center}

\subsection{Tipo de Conexões}

\begin{itemize}
    \item \textbf{Conexões duradouras:} entre peer e seu tracker. O peer mantém a conexão aberta para enviar notificações (ex.: Eu quero fazer o upload dessa musica).
    \item \textbf{Conexões temporárias:} entre peers durante download/upload. São estabelecidas sob demanda, após a descoberta dos peers fonte, e encerradas ao término da transferência.
\end{itemize}

% ============================================================
\section{Protocolos e Mensagens}\label{sec:mensagens}
% ============================================================
\subsection{Formato}

Todas as mensagens de controle (registro, busca, sincronização) utilizam \textbf{JSON} sobre HTTP/TCP, por ser legível e amplamente suportado. 

A transferência de chunks de áudio ocorre em \textbf{binário puro} sobre TCP, sem encapsulamento JSON, para não impor overhead desnecessário em dados já comprimidos.

\subsection{Tipos de Mensagem}

\begin{lstlisting}[language=Python, caption={Definição dos tipos de mensagem}, label=lst:mensagens]
# =====================================================================
# 1. Entrada do peer na rede e manutencao de presenca
# =====================================================================

# Apresentacao inicial do peer ao tracker (peer -> tracker)
# E a primeira mensagem que um peer envia ao se conectar; estabelece
# sua presenca antes de qualquer operacao de upload, busca ou download.
PEER_HELLO = {
    "type":      "PEER_HELLO",
    "nome_peer": "string",
    "ip":        "string",
    "porta":     "int"
}

# Peer saindo da rede de forma ordenada (peer -> tracker)
PEER_LEAVE = {
    "type":      "PEER_LEAVE",
    "nome_peer": "string"
}

# Notificacao de mudanca de IP (peer -> tracker)
# O tracker que recebe propaga a mudanca para os demais via SYNC_TABLE.
UPDATE_IP = {
    "type":      "UPDATE_IP",
    "nome_peer": "string",
    "novo_ip":   "string",
    "porta":     "int"
}

# Relatorio periodico de seed (peer -> tracker, a cada 3 minutos)
# Funciona tanto como anti-entropy do indice quanto como sinal de vida
# (failure detection): peer ausente por 2 rodadas e considerado falho.
SEED_REPORT = {
    "type":      "SEED_REPORT",
    "nome_peer": "string",
    "ip":        "string",
    "porta":     "int",
    "hashes":    ["sha256_hex"]
}

# =====================================================================
# 2. Registro de arquivos e busca
# =====================================================================

# Registro de musica (peer -> tracker)
# Usado em dois momentos: (a) upload original, com todos os campos;
# (b) ao concluir um download (fluxo de busca e download),
# o peer reenvia REGISTER_FILE para passar a constar como nova fonte.
# Nesse caso os campos nome/tamanho/n_chunks sao opcionais, o tracker
# ja os conhece a partir do upload original.
REGISTER_FILE = {
    "type":      "REGISTER_FILE",
    "nome_peer": "string",
    "hash":      "sha256_hex",
    "nome":      "string  (opcional ao re-registrar)",
    "tamanho":   "int_bytes (opcional ao re-registrar)",
    "n_chunks":  "int    (opcional ao re-registrar)"
}

# Busca de musica (peer -> tracker)
# O query_id e gerado pelo peer e correlaciona pedido e resposta,
# permitindo buscas paralelas. O ttl limita o roteamento entre trackers.
SEARCH_FILE = {
    "type":     "SEARCH_FILE",
    "query_id": "uuid",
    "query":    "string",
    "ttl":      "int"   # valor inicial recomendado: 3
}

# Roteamento de busca entre trackers (tracker -> tracker)
# Encaminhada quando o tracker local nao encontra o nome procurado.
# Carrega o tracker de origem para que a resposta volte direto.
SEARCH_FORWARD = {
    "type":           "SEARCH_FORWARD",
    "query_id":       "uuid",
    "query":          "string",
    "ttl":            "int",        # decrementado a cada salto
    "origem_tracker": "tracker_id"  # para onde devolver SEARCH_RESULT
}

# Resposta com peers disponiveis (tracker -> peer  ou  tracker -> tracker)
# Lista vazia em 'resultados' significa "nada encontrado".
# Cada resultado carrega 'n_chunks' (conhecido do REGISTER_FILE original)
# para o peer montar o plano de download sem precisar inferir o total a
# partir dos CHUNK_LIST das fontes, o que falharia se nenhuma fonte
# tivesse o arquivo completo.
SEARCH_RESULT = {
    "type":       "SEARCH_RESULT",
    "query_id":   "uuid",
    "resultados": [
        {
            "hash":     "sha256_hex",
            "nome":     "string",
            "n_chunks": "int",
            "peers": [
                { "nome_peer": "string", "ip": "string", "porta": "int" }
            ]
        }
    ]
}

# =====================================================================
# 3. Transferencia peer-to-peer
# =====================================================================

# Lista de chunks que um peer possui (peer -> peer)
# Antes de baixar, o peer solicitante pergunta a cada fonte qual subset
# de chunks ela tem; essencial para downloads parciais e paralelos.
CHUNK_LIST_REQUEST = {
    "type": "CHUNK_LIST_REQUEST",
    "hash": "sha256_hex"
}

CHUNK_LIST = {
    "type":         "CHUNK_LIST",
    "hash":         "sha256_hex",
    "chunks_disponiveis": ["int"]   # indices dos chunks que esse peer tem
}

# Requisicao de chunk (peer -> peer)
CHUNK_REQUEST = {
    "type":        "CHUNK_REQUEST",
    "hash":        "sha256_hex",
    "chunk_index": "int"
}

# Resposta de chunk (peer -> peer)
# Enviada como header JSON + payload binario na mesma conexao TCP.
CHUNK_DATA_HEADER = {
    "type":          "CHUNK_DATA",
    "hash":          "sha256_hex",
    "chunk_index":   "int",
    "payload_bytes": "int"   # tamanho do binario que vem em seguida
}

# =====================================================================
# 4. Remocao de arquivos
# =====================================================================

# Remocao explicita de um arquivo do indice (peer -> tracker)
# Redundante com SEED_REPORT (omitir o hash no proximo relatorio teria
# o mesmo efeito), mas existe para resposta imediata, sem esperar o
# proximo ciclo de 3 minutos.
PEER_LEAVE_FILE = {
    "type":      "PEER_LEAVE_FILE",
    "nome_peer": "string",
    "hash":      "sha256_hex"
}

# =====================================================================
# 5. Sincronizacao entre trackers (flooding TCP unicast)
# =====================================================================

# Atualizacao incremental do indice (tracker -> tracker)
# Cada entry carrega o IP:porta do peer (necessario para o tracker
# destinatario responder a buscas) e a flag 'ativo' que distingue
# adicao de remocao (tombstone).
# Entries de registro (ativo: True) carregam tambem os metadados do
# arquivo (nome, tamanho, n_chunks), pelos mesmos motivos do FULL_SYNC:
# sem eles, um tracker que conhece o hash apenas via SYNC_TABLE nao
# consegue responder buscas por nome nem aceitar o re-registro
# pos-download de um peer local. Tombstones (ativo: False) os omitem.
SYNC_TABLE = {
    "type":      "SYNC_TABLE",
    "origem":    "tracker_id",
    "timestamp": 1743550800.123,    # Unix timestamp (NTP-sincronizado)
    "entries": [
        {
            "hash":      "sha256_hex",
            "nome_peer": "string",
            "ip":        "string",
            "porta":     "int",
            "ativo":     "bool",     # True = registro, False = tombstone
            "nome":      "string  (opcional; presente quando ativo=True)",
            "tamanho":   "int_bytes (opcional; presente quando ativo=True)",
            "n_chunks":  "int    (opcional; presente quando ativo=True)"
        }
    ]
}

# Sincronizacao completa do indice (tracker -> tracker)
# Enviada em dois casos: (a) em resposta a TRACKER_REJOIN e (b)
# periodicamente, como reconciliacao anti-entropy entre trackers. Carrega o
# estado completo das tabelas e a lista atual de trackers ativos. Cada peer
# leva 'origem' (tracker que produziu a escrita), alem de 'ativo'/'timestamp',
# para o receptor aplicar LWW de forma deterministica (mesmo desempate do
# SYNC_TABLE).
FULL_SYNC = {
    "type":   "FULL_SYNC",
    "origem": "tracker_id",
    "entries": [
        {
            "hash":      "sha256_hex",
            "nome":      "string",
            "tamanho":   "int_bytes",
            "n_chunks":  "int",
            "peers": [
                {
                    "nome_peer": "string", "ip": "string", "porta": "int",
                    "ativo": "bool", "timestamp": "float",
                    "origem": "tracker_id"   # tracker que produziu a escrita (desempate LWW)
                }
            ]
        }
    ],
    "trackers_conhecidos": [
        { "tracker_id": "string", "ip": "string", "porta": "int" }
    ]
}

# =====================================================================
# 6. Gestao de membros do grupo de trackers
# =====================================================================

# Novo tracker se apresentando ao bootstrap node (tracker -> tracker)
TRACKER_REJOIN = {
    "type":       "TRACKER_REJOIN",
    "tracker_id": "string",
    "ip":         "string",
    "porta":      "int"
}

# Anuncio de novo tracker aos demais (tracker -> tracker, flooding TCP)
# Emitido pelo bootstrap node apos receber TRACKER_REJOIN.
TRACKER_ANNOUNCE = {
    "type": "TRACKER_ANNOUNCE",
    "novo_tracker": {
        "tracker_id": "string",
        "ip":         "string",
        "porta":      "int"
    }
}

# Reatribuicao de peer ao tracker reintegrado (tracker ATIVO -> peer)
# Apesar do nome, esta mensagem e enviada AO PEER (nao a outro tracker):
# informa que o peer deve agora se reportar a outro tracker, e nao mais
# ao atual. Usada no rebalanceamento.
REASSIGN_TRACKER = {
    "type":               "REASSIGN_TRACKER",
    "peer_nome":          "string",
    "novo_tracker_ip":    "string",
    "novo_tracker_porta": "int"
}

# =====================================================================
# 7. Sinalizacao de erro generica
# =====================================================================

# Resposta de erro a qualquer requisicao (qualquer -> qualquer)
# Usada quando uma operacao nao pode ser completada: hash inexistente,
# chunk indisponivel, peer desconhecido, etc.
ERROR = {
    "type":      "ERROR",
    "ref_type":  "string",   # tipo da mensagem que falhou (ex: "CHUNK_REQUEST")
    "ref_id":    "string",   # query_id ou outro correlator, se aplicavel
    "code":      "string",   # ex: "NOT_FOUND", "INVALID_HASH", "UNAUTHORIZED"
    "mensagem":  "string"    # descricao legivel para depuracao
}
\end{lstlisting}

\subsubsection*{Correlação de buscas}

Os campos \texttt{query\_id} em \texttt{SEARCH\_FILE}, \texttt{SEARCH\_FORWARD} e \texttt{SEARCH\_RESULT} são UUIDs gerados pelo peer no momento da busca; permitem correlacionar requisição e resposta, possibilitando que o peer faça múltiplas buscas em paralelo sem confundir os resultados. O \texttt{SEARCH\_FORWARD} carrega ainda o \texttt{origem\_tracker} para que o tracker remoto envie a resposta diretamente ao tracker que originou o roteamento, sem refazer todo o caminho.

\subsubsection*{Sobre a utilidade do \texttt{SEARCH\_FORWARD} com índice replicado}

Como o índice é integralmente replicado entre os trackers (Seção~\ref{sec:replicacao}), em estado estacionário toda busca é respondida localmente e o \texttt{SEARCH\_FORWARD} é dispensável. Ele permanece no protocolo como rede de segurança para a janela transitória de inconsistência: tracker recém-reintegrado ainda sem \texttt{FULL\_SYNC} completo, ou que perdeu um \texttt{SYNC\_TABLE} e ainda não foi reconciliado pelo \textit{anti-entropy}.

\subsubsection*{Limite de roteamento}
O campo \texttt{ttl} em \texttt{SEARCH\_FILE} é decrementado a cada
\texttt{SEARCH\_FORWARD} entre trackers, evitando que uma busca percorra a rede indefinidamente. O valor inicial recomendado é 3, suficiente para a topologia de poucos trackers do projeto.

\subsubsection*{Transferência de chunks}
A resposta \texttt{CHUNK\_DATA} é enviada na mesma conexão TCP da requisição: primeiro o cabeçalho JSON (\texttt{CHUNK\_DATA\_HEADER}, mostrado acima) com o campo \texttt{payload\_bytes} indicando o tamanho, seguido imediatamente pelo payload binário do chunk em si. Esse formato evita o overhead de codificar dados binários (já comprimidos como áudio) em JSON. Antes de iniciar a transferência, o peer cliente envia \texttt{CHUNK\_LIST\_REQUEST} a cada fonte para descobrir quais chunks específicos cada uma possui, isso permite paralelizar o download entre fontes que tenham chunks complementares.

\subsubsection*{Redundância intencional entre \texttt{PEER\_LEAVE\_FILE} e \texttt{SEED\_REPORT}}
Em estado estacionário, o \texttt{SEED\_REPORT} periódico já mantém o índice consistente: hashes ausentes no relatório acabam virando tombstones. O \texttt{PEER\_LEAVE\_FILE} é mantido como caminho de \textit{resposta imediata}: quando o usuário remove uma música, faz sentido o tracker propagar a remoção em segundos em vez de aguardar até três minutos pelo próximo \texttt{SEED\_REPORT}.

\subsubsection*{Sinalização de erro}
A mensagem genérica \texttt{ERROR} é usada como resposta negativa a qualquer requisição que não possa ser atendida, por exemplo, um \texttt{CHUNK\_REQUEST} para um chunk que o peer fonte já removeu, ou um \texttt{REGISTER\_FILE} com hash inválido. Os campos \texttt{ref\_type} e \texttt{ref\_id} permitem ao remetente correlacionar o erro com a requisição original.

\subsection{Fluxos Principais}\label{sec:fluxos}

Definidos os tipos de mensagem, descrevemos abaixo os fluxos centrais do sistema.

\subsubsection*{Entrada do peer na rede}
\begin{enumerate}
    \item Peer estabelece conexão TCP duradoura com seu tracker.
    \item Peer envia \texttt{PEER\_HELLO(nome\_peer, ip, porta)} para se registrar como presente.
    \item Tracker insere o peer em sua tabela \texttt{nome\_peer $\rightarrow$ (ip, porta, last\_seed\_ts)}; nenhum arquivo é registrado aqui.
\end{enumerate}

\subsubsection*{Upload de música}
\begin{enumerate}
    \item Pré-requisito: o peer já enviou \texttt{PEER\_HELLO} (sem isso o tracker rejeita o registro com \texttt{ERROR/PEER\_UNKNOWN}).
    \item Peer faz hash SHA-256 do arquivo localmente.
    \item Peer envia \texttt{REGISTER\_FILE(nome\_peer, hash, nome, tamanho, n\_chunks)} ao seu tracker via REST.
    \item Tracker atualiza as tabelas locais: \texttt{nome→hashes}, \texttt{hash→metadados}, \texttt{hash→peers}.
    \item Tracker envia \texttt{SYNC\_TABLE(origem, timestamp, entries[\{hash, nome\_peer, ip, porta, ativo=True, nome, tamanho, n\_chunks\}])} via TCP unicast para cada tracker conhecido, em paralelo.
    \item Cada tracker receptor aplica a entrada por LWW (Seção~\ref{sec:rep-protocolo}).
\end{enumerate}

\subsubsection*{Busca e download}
\begin{enumerate}
    \item Peer envia \texttt{SEARCH\_FILE(query\_id, query, ttl=3)} ao seu tracker via REST.
    \item Tracker consulta \texttt{nome→hashes} e, para cada hash, monta \texttt{hash→\{metadados, peers\}}.
    \item Se nada encontrado e \texttt{ttl > 0}, encaminha \texttt{SEARCH\_FORWARD} aos demais trackers em paralelo (Seção~\ref{sec:rep-protocolo}); agrega as respostas recebidas.
    \item Tracker retorna \texttt{SEARCH\_RESULT(query\_id, resultados=[\{hash, nome, n\_chunks, peers[\{nome\_peer, ip, porta\}]\}])} ao peer.
    \item Peer envia \texttt{CHUNK\_LIST\_REQUEST(hash)} a cada peer fonte; cada um responde com \texttt{CHUNK\_LIST(hash, chunks\_disponiveis)}.
    \item Peer monta um plano de download paralelo, distribuindo os chunks entre as fontes disponíveis, e envia \texttt{CHUNK\_REQUEST(hash, chunk\_index)} para cada chunk. Cada fonte responde com \texttt{CHUNK\_DATA\_HEADER(payload\_bytes)} seguido do payload binário, na mesma conexão TCP.
    \item Após completar o download e validar o SHA-256 do arquivo final, o peer envia \texttt{REGISTER\_FILE(nome\_peer, hash)} ao tracker (re-registro: \texttt{nome}/\texttt{tamanho}/\texttt{n\_chunks} são opcionais), passando a ser também fonte do arquivo.
\end{enumerate}

O último passo do download faz de cada peer que conclui uma transferência uma nova fonte do arquivo.

\subsubsection*{Remoção de música}
\begin{enumerate}
    \item Peer envia \texttt{PEER\_LEAVE\_FILE(nome\_peer, hash)} ao seu tracker via REST (ou \texttt{PEER\_LEAVE} para sair da rede inteira).
    \item Tracker remove o peer da lista de fontes do hash e cria um \textit{tombstone} local com \textit{timestamp} atual.
    \item Tracker envia \texttt{SYNC\_TABLE(origem, timestamp, entries[\{hash, nome\_peer, ip, porta, ativo=False\}])} aos demais trackers em paralelo.
    \item Cada tracker receptor aplica o \textit{tombstone} por LWW; o registro é retido por 10 minutos antes do descarte definitivo.
\end{enumerate}

% ============================================================
\section{Diagramas de Sequência}\label{sec:diag-seq}
% ============================================================

Os dois diagramas a seguir detalham, no tempo, duas operações centrais já descritas em alto nível na Seção~\ref{sec:fluxos}: o upload de uma música (com a consequente sincronização entre trackers) e a busca seguida de download paralelo entre peers.

\subsection{Upload de Música}
\begin{figure}[H]
    \centering
    \includegraphics[width=1\linewidth]{upload_p2p.jpg}
    \caption{Diagrama de sequência do upload}
    \label{fig:seq-upload}
\end{figure}

\subsection{Busca e Download de Música}
\begin{figure}[H]
    \centering
    \includegraphics[width=1\linewidth]{download_p2p.jpg}
    \caption{Diagrama de sequência do download}
    \label{fig:seq-download}
\end{figure}

% ============================================================
\section{Nomeação}\label{sec:nomeacao}
% ============================================================

Tanenbaum, no Cap. 6.1, estabelece que um \textit{true identifier} satisfaz três propriedades: 
\begin{enumerate}
    \item refere-se a no máximo uma entidade
    \item cada entidade é referida por no máximo um identificador
    \item nunca é reutilizado
\end{enumerate}

 No PeerSpot, três categorias de entidade precisam ser nomeadas: peers, arquivos de áudio e trackers.

\subsubsection*{Peers}
Cada nó da rede precisa de um identificador único para que o tracker saiba quem possui quais arquivos e para que peers se conectem diretamente entre si durante a transferência.
    
\subsubsection*{Arquivos de áudio}
As músicas precisam ser identificadas de maneira única e independente de nome, dois arquivos com o mesmo nome mas conteúdo diferentes são entidades diferentes.
    
\subsubsection*{Trackers}
Precisam se identificar entre si para controlar a origem das mensagens \texttt{SYNC\_TABLE} e evitar processamento de atualizações duplicadas.


\subsection{Esquema de Nomeação}\label{sec:esquema-nome}

O PeerSpot adota esquemas complementares, conforme a entidade:

\subsubsection*{Nomes estruturados para peers}

Peers são identificados por um \textbf{nome de usuário}, que deve ser único na rede. O tracker mantém uma tabela \texttt{nome\_peer $\rightarrow$ ip:porta} que mapeia cada peer ao seu endereço de rede atual. Como peers domésticos podem ter IPs dinâmicos, toda vez que o IP de um peer muda ele notifica seu tracker com uma mensagem \texttt{UPDATE\_IP}, que propaga a atualização aos demais trackers via flooding TCP. Dessa forma, o nome do peer permanece estável como identidade, enquanto o IP é tratado como endereço de contato mutável.

O \texttt{nome\_peer} satisfaz as propriedades 1. e 2. de um \textit{true identifier}, mas viola estritamente a 3.: nada impede que um usuário desligue sua conta e outro mais tarde se cadastre com o mesmo nome. Para o escopo do projeto isso é aceitável, visto que não há reuso esperado em janelas de tempo curtas.

\subsubsection*{Nomes planos para arquivos}

Arquivos de áudio são identificados pelo \textbf{hash SHA-256} do conteúdo. Esse é um nome plano, uma sequência de 256 bits sem estrutura legível. 

Como é derivado do conteúdo, satisfaz integralmente as três propriedades de um \textit{true identifier}: dois arquivos com mesmo conteúdo recebem o mesmo identificador (o que é desejável, pois são a mesma entidade lógica) e conteúdos diferentes nunca colidem.

\subsubsection*{Identificação dos trackers}

Cada tracker recebe um \texttt{tracker\_id} único, atribuído administrativamente no momento da configuração da rede (ex.: \texttt{tracker-1}, \texttt{tracker-2}, ...). Esse identificador aparece no campo \texttt{origem} de toda mensagem \texttt{SYNC\_TABLE} e serve dois propósitos: evitar reprocessamento de mensagens recebidas como eco do próprio envio, e desempatar conflitos de \textit{Last Write Wins} com timestamps idênticos. Como o conjunto de trackers é gerenciado por configuração e não cresce dinamicamente em larga escala, a unicidade do \texttt{tracker\_id} é garantida operacionalmente.

\subsubsection*{Nomes estruturados das músicas}

Para a interface do usuário, o sistema mantém no tracker um mapeamento de nome legível para hash: \texttt{nome\_da\_musica} $\rightarrow$ \texttt{lista de hashes}.

Note que uma busca por nome pode retornar \textbf{múltiplos hashes}, pois é possível haver mais de uma música com o mesmo nome (artistas diferentes, versões distintas). Cabe ao usuário selecionar qual resultado deseja baixar. Internamente, toda operação usa o hash, o nome legível é apenas a camada de apresentação.

\subsection{Mecanismo de Resolução de Nomes}

A resolução envolve \textbf{duas tabelas consultadas no tracker, mas uma única troca de mensagens} entre peer e tracker. Ao receber um \texttt{SEARCH\_FILE(nome)}, o tracker:

\begin{enumerate}
    \item Consulta a tabela \texttt{nome $\rightarrow$ hashes} para obter todos os identificadores correspondentes ao nome buscado (suporta arquivos homônimos).
    \item Para cada hash, consulta as tabelas \texttt{hash $\rightarrow$ metadados} e \texttt{hash $\rightarrow$ peers} para anexar \texttt{n\_chunks} e a lista de fontes ativas.
    \item Devolve um único \texttt{SEARCH\_RESULT} com todos os resultados agregados, no formato \texttt{resultados=[\{hash, nome, n\_chunks, peers[\{nome\_peer, ip, porta\}]\}]}.
\end{enumerate}

Cabe ao peer (e ao usuário, na CLI) escolher entre os resultados retornados. Não há uma segunda chamada \texttt{hash $\rightarrow$ peers} sobre a rede: as duas tabelas são consultadas internamente pelo tracker antes da resposta.


Se o tracker local não possui o hash na tabela (\textit{race} com a propagação por flooding ou tracker recém-reintegrado), ele roteia a consulta para os demais trackers via TCP, com decremento de TTL para evitar loops. A sequência completa de mensagens é:

\begin{lstlisting}[basicstyle=\ttfamily\footnotesize, frame=single, breaklines=true]
peer -> tracker:  SEARCH_FILE(query_id, query, ttl=3)
  |
  +-- hit local --> SEARCH_RESULT(query_id,
                       [{hash, nome, n_chunks, peers[]}])
  |
  +-- sem hit e ttl>0:
        tracker -> trackers (paralelo):
            SEARCH_FORWARD(query_id, query, ttl-1, origem_tracker)
        cada receptor responde na MESMA conexao TCP:
            SEARCH_RESULT(query_id, [...])
        agrega respostas -> SEARCH_RESULT ao peer

peer <-> fontes:  CHUNK_LIST_REQUEST(hash) -> CHUNK_LIST(hash, chunks_disponiveis)
                  CHUNK_REQUEST(hash, chunk_index) -> CHUNK_DATA_HEADER + payload

peer -> tracker:  REGISTER_FILE(nome_peer, hash)   # re-registro como nova fonte
\end{lstlisting}

Vale destacar duas decisões: o \texttt{SEARCH\_RESULT} de um forward volta na mesma conexão TCP do encaminhamento (sem precisar abrir nova conexão de volta), e o tracker receptor de um \texttt{SEARCH\_FORWARD} não re-encaminha --- evita loops sem precisar de tabela de \textit{routing} adicional.

% ============================================================
\section{Processos}\label{sec:processos}
% ============================================================

A seguir analisamos cada aspecto para o PeerSpot.

\subsection{Faz sentido usar threads?}

Sim, e em ambos os componentes.

\subsubsection*{No tracker}
O tracker precisa atender múltiplos peers simultaneamente: receber registros de upload, responder buscas e processar mensagens de sincronização vindas de outros trackers. Um modelo single-threaded bloquearia todas as operações enquanto uma requisição estivesse sendo processada. 

A solução é um servidor multithread com pool de threads, onde cada requisição recebida é despachada para uma thread do pool roda em uma thread dedicada: um servidor TCP \texttt{accept()} ininterrupto aguarda conexões dos demais trackers e despacha cada uma para uma thread filha, em paralelo com o atendimento de requisições REST dos peers locais.

\subsubsection*{Exclusão mútua local entre as threads do tracker}
Como várias threads do tracker podem tocar as mesmas tabelas de índice em memória ao mesmo tempo (ex: uma atendendo um \texttt{REGISTER\_FILE} de um peer local enquanto outra aplica um \texttt{SYNC\_TABLE} recebido de outro tracker), o acesso a essas estruturas é protegido por \textbf{exclusão mútua local}, com um \texttt{threading.Lock()} padrão da linguagem. 

Trata-se de um problema de concorrência dentro de um processo em um único tracker, e não de exclusão mútua \textit{distribuída} entre trackers. A semântica comutativa das operações sobre o índice (somada ao LWW para remoções) dispensa qualquer coordenação de exclusão mútua entre os trackers.

\subsubsection*{No peer}
O peer precisa, simultaneamente, manter a conexão duradoura com o tracker, servir chunks para outros peers que estão baixando, baixar chunks de outras fontes, e manter a interface do usuário responsiva. Essas quatro atividades são naturalmente concorrentes e se beneficiam de threads independentes para cada responsabilidade.

\subsection{Servidores stateful ou stateless?}

No PeerSpot, consideramos o tracker como stateful e o peer como parcialmente stateful.

\subsubsection*{Tracker}

O tracker é stateful: mantém em memória as tabelas \texttt{nome $\rightarrow$ hashes}, \texttt{hash $\rightarrow$ peers} e \texttt{nome\_peer $\rightarrow$ ip:porta}. O estado é essencial para a função do tracker, visto que nenhuma busca pode ser respondida sem ele.

A implicação é que, ao reiniciar, o tracker começa com as tabelas vazias e reconstrói o índice via \texttt{FULL\_SYNC} com os demais trackers ativos.

\subsubsection*{Peer}

O peer é stateful no que diz respeito aos chunks armazenados localmente (persistidos em disco). A conexão duradoura com o tracker também é estado mantido em memória. 

Porém, as conexões temporárias de transferência entre peers são tratadas de forma stateless, dado que cada \texttt{CHUNK\_REQUEST} é independente, e a conexão TCP pode ser encerrada e reestabelecida sem perda de progresso, pois o gerenciador de chunks controla quais índices já foram recebidos.

\subsection{Faz sentido usar técnicas de virtualização?}

Faria sentido em um cenário de implantação real. A ideia seria usar containers, que são mais leves que uma máquina virtual porque compartilham o kernel do hospedeiro e ainda empacotam todo o ambiente de execução (bibliotecas, dependências) junto com a aplicação. Isso ajudaria com a \textbf{heterogeneidade} de um jeito diferente do middleware: em vez de abstrair as diferenças entre os sistemas, o container simplesmente as elimina, já que cada peer rodaria com o mesmo ambiente empacotado.

\paragraph{Em implantação real.}
Em um cenário de implantação, os trackers poderiam rodar em máquinas virtuais ou containers em provedores de nuvem, permitindo escalar o número de trackers conforme a demanda e recriar instâncias automaticamente em caso de falha. Entretanto, dado o escopo do nosso projeto, tamanho trabalho não é viável por questões de tempo de implementação.

\paragraph{Durante o desenvolvimento e testes (nossa escolha).}
Para o escopo deste projeto, optamos por uma abordagem mais simples e direta: simular múltiplos peers e trackers como \textbf{processos independentes em terminais separados} na mesma máquina, distinguidos por porta (\texttt{ip:porta}). Cada terminal executa um processo Python, um peer ou um tracker, com seu próprio arquivo de configuração. Como o modelo de endereçamento da rede já é baseado em \texttt{ip:porta}, rodar todos os nós em \texttt{127.0.0.1} com portas distintas reproduz fielmente a topologia distribuída sem nenhuma camada de virtualização.

Essa escolha permite testar os mesmos cenários que conteinerização permitiria, falha de nó (encerrar o processo com \texttt{Ctrl+C} ou \texttt{kill}), sincronização entre trackers, reintegração via \texttt{TRACKER\_REJOIN} e download paralelo entre peers, com custo de configuração muito menor e sem dependências externas além do interpretador Python. A heterogeneidade de ambiente, que a conteinerização resolveria, não é uma preocupação aqui, pois todos os nós rodam o mesmo código na mesma máquina durante o desenvolvimento.

% ============================================================
\section{Replicação e Consistência}\label{sec:replicacao}
% ============================================================

A seguir, descrevemos o que replicamos, sob qual modelo de consistência e por qual protocolo as cópias são mantidas em sincronia.

\subsection{O que é replicado}\label{sec:rep-oque}

A replicação acontece sobre os indíces e arquivos.

\subsubsection*{Índice (metadados)} 
Todos os trackers formam uma rede totalmente conectada e mantêm cópias completas do índice global (\textit{full replication}). 

Cada tracker é uma réplica integral. A motivação é tolerância a falhas: a queda de um tracker não derruba a busca e desempenho de leitura (toda busca é respondida localmente).

\subsubsection*{Arquivos (conteúdo)} 
O áudio é replicado naturalmente entre os peers: quando um peer termina um download, ele se registra como nova fonte do arquivo. Cada cópia baixada aumenta a disponibilidade e a banda total;

Como os arquivos são endereçados por conteúdo (hash SHA-256), eles são imutáveis, ou seja, duas réplicas de um mesmo hash são idênticas e nunca entram em conflito. Por isso, todo o problema de consistência recai apenas sobre o índice.

\subsection{Modelo de consistência}\label{sec:modelo_consistencia}

O índice adota \textbf{consistência eventual}: aceitamos pequenas inconsistências transitórias (no máximo alguns minutos) em troca de uma implementação simples e sempre disponível. Na ausência de novas atualizações, todas as réplicas convergem para o mesmo estado. 

A consistência forte foi descartada porque exigiria coordenação síncrona entre trackers a cada escrita, contrariando os objetivos de disponibilidade e simplicidade. No pior caso, uma busca retorna um peer que acabou de sair, e o cliente apenas tenta a próxima fonte.

Essa escolha se apoia na semântica das operações. Registrar um peer como fonte de um arquivo é uma operação aditiva e comutativa: dois trackers adicionando peers diferentes ao mesmo arquivo, em qualquer ordem, chegam ao mesmo resultado. 

Isso elimina conflitos no caso dominante. Já a remoção de um peer não é comutativa com a adição, e por isso é tratada com \textit{tombstones} e resolvida por \textit{Last Write Wins} (descritos adiante). Como consequência, a única exclusão mútua necessária é local, resolvida com locks de threading (Seção~\ref{sec:processos}).

\subsubsection*{Relógio real com NTP}
Para ordenar remoções conflitantes, usamos timestamps de relógio real, assumindo que os trackers têm NTP sincronizado. Como as atualizações de índice têm granularidade de segundos, a precisão do NTP é mais que suficiente. 

Relógios lógicos de Lamport ofereceriam ordenação causal estrita, que não é necessária aqui, já que a operação dominante é comutativa e não há dependência causal a preservar.

\subsubsection*{Resolução de conflitos: Last Write Wins}
Quando um tracker recebe uma atualização para uma entrada que já possui, compara os timestamps: se o recebido for mais recente, sobrescreve a versão local; caso contrário, descarta. 

É uma política simples e adequada ao domínio, já que o estado de um arquivo (quais peers o possuem) é naturalmente substituível pela versão mais nova.

\subsubsection*{Remoção com tombstones}
Quando um peer sai ou remove um arquivo, em vez de apagar a entrada na hora, o tracker a marca como removida por um período. 

Isso evita o seguinte problema: se um tracker estava desconectado e perdeu a mensagem de remoção, ele continuaria servindo um peer que já saiu. Com o tombstone, ao voltar à rede o tracker aplica o LWW normalmente e efetiva a remoção. O período de retenção é de \textbf{10 minutos}, escolhido para cobrir partições curtas e o ciclo de reconciliação \textit{anti-entropy} entre trackers (Seção~\ref{sec:rep-protocolo}); passado esse prazo, o registro é descartado de vez.

\subsubsection*{Distribuição das cópias}

O conjunto de trackers é \textbf{estático} (definido na configuração da rede). Já as réplicas dos arquivos são \textbf{dinâmicas e orientadas a demanda}: uma nova cópia surge exatamente quando e onde um peer baixa o arquivo.

\subsection{Protocolo de consistência}\label{sec:rep-protocolo}

O índice é mantido por escrita replicada sem primário: qualquer tracker aceita uma escrita e a propaga aos demais, sem coordenador. O protocolo combina propagação por \textit{flooding} TCP, resolução por LWW, remoções por tombstone e reconciliação \textit{anti-entropy}.

\subsubsection*{Propagação por flooding sobre TCP}
Cada tracker mantém a lista dos demais e, ao atualizar seu índice, envia a atualização por TCP unicast para todos em paralelo. Optamos por TCP em vez de IP multicast sobre UDP por dois motivos: 

\begin{enumerate}
    \item UDP não é confiável e exigiria mecanismos extras de retransmissão
    \item Para o número pequeno de trackers (tipicamente 2--5), abrir algumas conexões TCP é trivial e já dá entrega confiável, ordenada e detecção de falha pelo timeout
\end{enumerate}

Como a topologia é totalmente conectada, o \textit{flooding} aqui é ideal, já que cada mensagem vai uma única vez a cada destino, sem retransmissão, loops ou necessidade de TTL.

\subsubsection*{Entrada dinâmica de trackers}
Um novo tracker é iniciado com o endereço de pelo menos um tracker já ativo (\textit{bootstrap node}). Ele se anuncia a esse nó, recebe de volta o estado completo do índice e a lista de trackers ativos, e o \textit{bootstrap} divulga sua existência aos demais. A partir daí o novo tracker está integrado e passa a receber as atualizações.

\subsubsection*{Reconciliação anti-entropy}

O \texttt{SYNC\_TABLE} é \textbf{incremental} (carrega só o que mudou) e enviado em modo \textit{best-effort}: o remetente não retransmite em caso de falha. Isso levanta uma pergunta legítima: se um tracker perder uma única mensagem \texttt{SYNC\_TABLE} por uma falha transitória de TCP, mas continuar no ar, esse delta some do caminho de \textit{flooding} e ele ficaria desatualizado indefinidamente. O \texttt{SYNC\_TABLE} sozinho, portanto, \textbf{não} fecha o problema de entropia.

A convergência é garantida por um segundo mecanismo, executado em segundo plano: a \textbf{reconciliação anti-entropy periódica entre trackers}. A cada \textbf{3 minutos}, cada tracker envia aos demais uma mensagem \texttt{FULL\_SYNC} com o estado completo do seu índice --- fontes ativas e \textit{tombstones}, cada qual com seu \textit{timestamp} e tracker de origem. O receptor aplica cada entrada pela mesma regra de \textit{Last Write Wins} do \texttt{SYNC\_TABLE}; como essa aplicação é idempotente e comutativa, reaplicar o estado completo não causa efeito colateral e apenas \textbf{repara o que estiver divergente}, qualquer que tenha sido a causa (mensagem perdida, partição curta ou tracker que reiniciou). O pior caso de desatualização passa a ser limitado a um intervalo de reconciliação, em vez de permanente.

É por isso que, quando um envio de \texttt{SYNC\_TABLE} falha, o remetente apenas marca o destino como suspeito e \textbf{não} reenvia na hora: a próxima rodada de anti-entropy corrige a diferença, sem precisar de filas de retransmissão por destino nem de confirmação (\textit{ACK}) por mensagem. Preferimos o anti-entropy periódico a um \texttt{SYNC\_TABLE} confiável com \textit{ACK} e retransmissão porque ele cobre, com um único mecanismo simples, também os casos que o \textit{ACK} não cobriria --- como a queda do próprio remetente antes de concluir a entrega, ou um destino fora do ar por muito tempo. O período de reconciliação é mantido bem menor que a retenção dos \textit{tombstones} (Seção~\ref{sec:modelo_consistencia}), para que uma remoção perdida seja reposta antes de o \textit{tombstone} expirar.

Em complemento, no sentido peer$\rightarrow$tracker, o \texttt{SEED\_REPORT} periódico (Seção~\ref{sec:mensagens}) reconcilia o índice \textit{local} do tracker ao qual o peer se reporta: hashes ausentes no relatório viram \textit{tombstone}, novos são registrados. E um tracker que tenha ficado muito tempo fora, ao voltar, recupera todo o índice de uma réplica viva via \texttt{TRACKER\_REJOIN}/\texttt{FULL\_SYNC} (Seção~\ref{sec:reintegracao}).

Protocolos baseados em primário foram descartados para não precisar eleger um tracker líder, e protocolos por quórum também, já que votação só se justifica quando há conflitos de escrita a serializar, o que a comutatividade das adições e o LWW das remoções tornam desnecessário.

\subsection{Limitações assumidas}\label{sec:limit_repl}

Por transparência, duas limitações conscientes:

\begin{enumerate}
    \item Como a replicação de conteúdo é puramente orientada a demanda, não há garantia de um número mínimo de cópias por arquivo, então um arquivo fica indisponível se seu único peer sai da rede
    \item Manter o índice inteiro em cada tracker é ótimo para latência e tolerância a falhas, mas tem custo de memória proporcional ao catálogo, o que só seria problemático em uma escala muito maior que a do projeto.
\end{enumerate}


% ============================================================
\section{Tolerância a Falhas}\label{sec:tolerancia}
% ============================================================

A seguir, detalhamos como o PeerSpot lida com as possíveis falhas do sistema.

\subsection{Não há tracker líder, não há necessidade de eleição}\label{sec:sem_lider}

No PeerSpot, todos os trackers são equivalentes e mantêm cópias sincronizadas do índice global. Não existe tracker primário nem hierarquia entre eles, portanto não há necessidade de algoritmo de eleição. 

Tanenbaum (Cap. 5.4) aponta que algoritmos de eleição são necessários apenas quando o sistema depende de um coordenador único, o que não é o caso aqui.

\subsection{Fallback no peer}\label{sec:fallback_peer}

Cada peer mantém localmente uma lista ordenada de trackers conhecidos: o tracker principal e um ou mais trackers de backup. Quando o peer detecta que o tracker principal não responde (timeout na conexão duradoura), tenta o próximo da lista automaticamente, sem qualquer coordenação com outros nós da rede.

A reorganização é imediata e local, o peer simplesmente passa a operar com o próximo tracker da lista, que já possui o índice completo sincronizado.

\subsection{Reintegração do tracker e balanceamento de carga}\label{sec:reintegracao}

Quando um tracker que estava fora do ar volta a funcionar, ele precisa: (1)~recuperar o estado atual do índice e (2)~receber de volta parte dos peers para distribuir a carga entre os trackers ativos.

Ao iniciar, o tracker envia uma mensagem \texttt{TRACKER\_REJOIN} ao \textit{bootstrap node} que responde com \texttt{FULL\_SYNC} contendo o estado completo das tabelas e a lista atual de trackers ativos. O \textit{bootstrap node} então propaga \texttt{TRACKER\_ANNOUNCE} aos demais trackers, que adicionam o tracker reintegrado às suas listas \texttt{PEER\_TRACKERS}. Em seguida, cada tracker ativo cede uma fração de seus peers ao tracker reintegrado, enviando a cada peer cedido uma mensagem \texttt{REASSIGN\_TRACKER} com o endereço do novo tracker.

O critério de balanceamento é simples: cada tracker ativo divide seu número atual de peers pelo número total de trackers na rede (incluindo o reintegrado) e cede o excedente. Isso distribui a carga de forma aproximadamente uniforme sem necessidade de coordenação centralizada.

\subsubsection*{Sobre a consistência do balanceamento}
Como cada tracker ativo calcula sua cessão de forma independente, há risco de \textit{sobrecessão}, todos os trackers cederem peers ao mesmo tempo e o tracker reintegrado acabar com mais peers que os demais. Esse desbalanceamento transitório é aceitável: na próxima rodada de \texttt{SEED\_REPORT}, cada tracker recalcula seu fator de carga e, se algum estiver substancialmente abaixo da média, novos \texttt{REASSIGN\_TRACKER} são emitidos no sentido inverso. O sistema converge para um balanceamento aproximadamente uniforme sem necessidade de coordenação atômica entre os trackers.

\subsubsection*{Sobre a lista \texttt{TRACKERS} no peer}
A lista de trackers que cada peer mantém localmente é estática no código apresentado. Para suportar a adição dinâmica de trackers, a lista pode ser estendida por uma resposta opcional do tracker em qualquer comunicação subsequente, incluindo um campo \texttt{trackers\_conhecidos} com o conjunto atual de trackers ativos. O peer atualiza sua lista local com essa informação. Para simplicidade do protótipo, no escopo deste trabalho assumimos a lista estática, mas a limitação é registrada aqui explicitamente.

\subsection{Disponibilidade ou Confiabilidade?}\label{sec:disp_conf}

Entre disponibilidade (estar pronto para uso a qualquer instante) e confiabilidade (operar continuamente, sem falhas, por longos períodos), o PeerSpot prioriza a \textbf{disponibilidade}. A razão vem do próprio domínio: em um serviço de compartilhamento de música, é melhor que uma busca responda rápido com um resultado talvez desatualizado (por exemplo, um peer que acabou de sair) do que travar a resposta esperando ter certeza absoluta.

Essa preferência é coerente com as escolhas que já fizemos: o índice é replicado em todos os trackers (qualquer um responde), o peer troca na hora para um tracker de backup quando o principal cai (Seção~\ref{sec:fallback_peer}) e o sistema opera sob consistência eventual (Seção~\ref{sec:modelo_consistencia}). Em troca, abrimos mão da confiabilidade estrita, aceitando pequenas inconsistências transitórias em favor da disponibilidade e da simplicidade.

\subsection{Tipos de Falha que se Deseja Tolerar}\label{sec:tipos-falha}

O PeerSpot assume o modelo \textit{crash} (\textit{fail-stop}): os nós podem parar, mas não passam a se comportar de forma arbitrária. Percorrendo os tipos de falha:

\begin{itemize}
    \item \textbf{Crash (parada):} \textbf{tolerada}. É o caso central. Peers e trackers podem parar de uma hora para outra, e a tolerância vem da replicação total do índice e da auto-replicação dos arquivos.
    \item \textbf{Omissão (perda de mensagens):} \textbf{mitigada}. Mensagens perdidas são repostas pela reconciliação \textit{anti-entropy} periódica entre trackers (\texttt{FULL\_SYNC}), pelo \texttt{SEED\_REPORT} no sentido peer$\rightarrow$tracker e pelo \texttt{FULL\_SYNC} na reintegração; o uso de TCP já evita a maior parte das perdas.
    \item \textbf{Falha de resposta (resposta incorreta):} \textbf{fora do escopo}. Assumimos que um nó que responde, responde certo. Na prática, uma resposta que não chega recai no caso de crash ou omissão.
    \item \textbf{Falha temporal:} \textbf{fora do escopo}. Assumimos o NTP funcionando, e pequenos desvios de relógio são absorvidos pelo LWW.
    \item \textbf{Arbitrária / bizantina:} \textbf{fora do escopo}. Não tratamos nós maliciosos ou com comportamento arbitrário; não há votação nem assinatura para detectar mentiras. Também não há autenticação de peers.
\end{itemize}

\subsection{Quantos Processos Falhantes Serão Suportados?}

Como o PeerSpot não usa votação por quórum nem acordo bizantino, não vale aqui o limite clássico de $k+1$ réplicas para tolerar $k$ falhas. A tolerância depende apenas do grau de replicação:

\begin{itemize}
    \item \textbf{Trackers:} como o índice está replicado por inteiro nos $N$ trackers, a busca continua funcionando enquanto pelo menos um deles estiver no ar, ou seja, toleramos a queda de até $N-1$ trackers ao mesmo tempo. Os peers migram sozinhos para o próximo tracker da lista (Seção~\ref{sec:fallback_peer}).
    \item \textbf{Arquivos:} um arquivo continua acessível enquanto pelo menos um peer que o possui estiver on-line. Se ele tem $r$ cópias na rede, toleramos a perda de até $r-1$ delas. Como o sistema não controla $r$ (a replicação é orientada a demanda), não há garantia de um número mínimo de cópias, limitação já registrada na Seção~\ref{sec:limit_repl}.
\end{itemize}

\subsection{Estratégias para Detectar Falhas}\label{sec:detec-falha}

A detecção é sempre baseada em \textbf{timeout}, seja sobre uma comunicação periódica, seja sobre uma conexão. O caso que realmente exige detecção ativa é a \textbf{saída desordenada} de um peer (queda de conexão, processo encerrado de repente, máquina desligada), em que não chega nenhum \texttt{PEER\_LEAVE} avisando da saída. No PeerSpot, há três caminhos de detecção:

\begin{itemize}
    \item \textbf{Falha de peer, detectada pelo tracker:} se um peer fica duas rodadas seguidas (6~min) sem enviar \texttt{SEED\_REPORT}, o tracker o considera falho, marca todas as suas entradas com \textit{tombstone} e propaga via \texttt{SYNC\_TABLE}. A detecção acaba sendo um efeito colateral da troca periódica de mensagens.
    \item \textbf{Falha de tracker, detectada pelo peer:} o \textit{timeout} na conexão duradoura dispara o fallback para o próximo tracker da lista (Seção~\ref{sec:fallback_peer}).
    \item \textbf{Falha de tracker, detectada por outro tracker:} um \textit{timeout} ou \texttt{ConnectionRefusedError} ao tentar enviar \texttt{SYNC\_TABLE} marca o tracker como suspeito.
\end{itemize}

Quando a saída é \textbf{ordenada}, não há o que detectar: o peer envia \texttt{PEER\_LEAVE} (ou \texttt{PEER\_LEAVE\_FILE}, para uma música específica) e o tracker propaga a remoção na hora.

\subsection{Qual Protocolo Utilizar?}

O PeerSpot não usa protocolo de consenso (Paxos, Raft) nem commit distribuído (2PC/3PC). Isso é uma decisão de projeto, não uma falta: como não há tracker líder (Seção~\ref{sec:sem_lider}) nem conflitos de escrita a serializar (Seção~\ref{sec:modelo_consistencia}), não existe nenhuma decisão sobre a qual os nós precisem chegar a um acordo atômico. A resiliência vem da combinação de quatro mecanismos:

\begin{enumerate}
    \item \textbf{Comunicação confiável ponto a ponto:} o TCP unicast garante entrega ordenada ou um erro explícito.
    \item \textbf{Detecção de falhas por timeout} (Seção~\ref{sec:detec-falha}).
    \item \textbf{Redundância por replicação:} índice replicado por inteiro e arquivos auto-replicados entre os peers.
    \item \textbf{Recuperação por transferência de estado:} o \texttt{FULL\_SYNC} reconstrói o índice de um tracker reintegrado (Seção~\ref{sec:reintegracao}).
\end{enumerate}

\subsection{Consequências do Teorema CAP}

O teorema CAP diz que, durante uma \textbf{partição de rede} (P), um sistema distribuído não consegue garantir ao mesmo tempo \textbf{consistência forte} (C) e \textbf{disponibilidade} (A): é preciso escolher uma das duas.

O PeerSpot é, por escolha, um sistema \textbf{AP}: diante de uma partição entre trackers, cada lado continua respondendo a buscas e aceitando registros com seu índice local (prioriza A), abrindo mão da consistência forte (sacrifica C) e operando sob consistência eventual. Quando a partição é curada, os índices voltam a convergir via \textit{anti-entropy} periódica entre trackers (\texttt{FULL\_SYNC}) e via \texttt{SEED\_REPORT}, com os conflitos resolvidos por LWW (Seção~\ref{sec:rep-protocolo}).

Na prática, isso significa que:
\begin{itemize}
    \item Durante uma partição, dois trackers podem ter visões momentaneamente diferentes do índice (um peer aparece de um lado e não do outro). Tudo bem: a inconsistência é passageira e o pior que acontece é uma busca falhar e precisar ser repetida.
    \item Em nenhum momento o sistema fica indisponível só para preservar consistência. A opção contrária (CP) exigiria, por exemplo, bloquear escritas sem quórum, o que iria contra o objetivo de disponibilidade da Seção~\ref{sec:disp_conf}.
    \item Como os arquivos são imutáveis (endereçados por hash), a partição nunca causa divergência de \textit{conteúdo}, só de \textit{disponibilidade de fontes}, o que deixa o custo de escolher A bem baixo.
\end{itemize}

\subsection{Como Recuperar da Falha?}

Os mecanismos de recuperação variam conforme o tipo de falha:

\begin{itemize}
    \item \textbf{Tracker que cai:} os peers migram na hora para um tracker de backup já sincronizado (Seção~\ref{sec:fallback_peer}), sem perder o serviço.
    \item \textbf{Tracker que retorna:} reconstrói todo o índice a partir de um tracker ativo, via \texttt{TRACKER\_REJOIN} seguido de \texttt{FULL\_SYNC} (Seção~\ref{sec:reintegracao}). É como restaurar o estado atual de uma réplica viva, sem precisar reaplicar um log de operações. Em seguida vem o rebalanceamento de carga.
    \item \textbf{Partição curta:} a reconciliação \textit{anti-entropy} periódica entre trackers (\texttt{FULL\_SYNC}) repõe sozinha as entradas que faltaram.
    \item \textbf{Peer que cai:} suas entradas viram \textit{tombstone} após o timeout do \texttt{SEED\_REPORT} (Seção~\ref{sec:detec-falha}); quando ele volta, basta se re-registrar para constar de novo como fonte.
\end{itemize}

Como o índice fica em memória e pode sempre ser reconstruído a partir das réplicas vivas ou dos próprios peers, os trackers não precisam de \textit{checkpoint} em disco nem de \textit{log} de mensagens: o ``checkpoint'', na prática, é o estado de qualquer réplica que esteja no ar.

\end{document}